\subsection{Objetivo de la compresión parcelaria}

Feature Table v1 transforma fuentes heterogéneas y longitudinales en una representación
determinista de una fila por parcela. El artefacto no es un dataset limpio en el sentido
de haber aplicado imputación global o escalado aprendido; es una tabla \emph{model-ready}
porque resuelve las uniones, unidades, agregaciones y provenance, dejando cualquier
transformación entrenable dentro de la validación.

El build validado produjo:
\begin{table}[H]
\centering
\caption{Dimensiones de Feature Table v1.}
\label{tab:cp02-featuretable}
\begin{tabular}{lrr}
\toprule
Artefacto & Filas & Columnas \\
\midrule
\artifact{parcel_features_clean.csv} & 197 & 699 \\
\artifact{parcel_features_competition.csv} & 197 & 1,408 \\
\artifact{parcel_features_all.csv} & 197 & 1,408 \\
\bottomrule
\end{tabular}
\end{table}

El manifest contiene 1,403 predictores: 694 clean y 709 competition-only. Las columnas
adicionales corresponden a ID/metadata/target/split. Todas las tablas preservan exactamente
197 IDs, 138 targets observados y 59 targets missing.

\subsection{Familias de variables}

El inventario de Checkpoint~02 clasifica las features por source, family y mode. Entre las
features clean aparecen 539 BASIC históricos, 90 SoilGrids, 24 de clima oficial histórico,
14 SIAP históricos, ocho de topografía, seis geométricas y cinco CEM, además de campos
administrativos. La capa competition añade 539 BASIC 2025, 97 PRO, 45 clima 2025, 22 WaPOR,
cinco SIAP contemporáneas y la feature CHIRPS disponible.

Esta taxonomía permite realizar ablaciones semánticas. No basta con preguntar si una columna
mejora RMSE; interesa identificar si el incremento proviene de geometría, clima, historial
satelital, señal contemporánea o una estadística administrativa.

\subsection{Agregación temporal}

Para una serie longitudinal $x_{i,s,t}$ de parcela $i$, sensor $s$ y fecha $t$, Feature Table
v1 genera resúmenes compatibles con el source contract: históricos, 2025, mensuales,
estacionales y por sensor. La operación conceptual es
\begin{equation}
x_{i,s,\mathcal T}^{(g)}
=
g\left(\{x_{i,s,t}:t\in\mathcal T\}\right),
\label{eq:feature-aggregation}
\end{equation}
donde $g$ puede ser media, desviación, mínimo, máximo u otro resumen explícito.

La agregación no cruza sensores por conveniencia. Por ejemplo, Sentinel-2 NDVI y Landsat NDVI
son variables distintas aun si comparten nombre del índice. Esta decisión evita que
disponibilidad temporal/sensor-specific missingness se convierta en un promedio sin
interpretación física.

La tabla final no pretende agotar el contenido temporal. El Checkpoint~04A vuelve a
trayectorias mensuales para construir correlación, lag y DTW. En otras palabras, Feature Table
v1 es una proyección tabular de los datos longitudinales, no una afirmación de suficiencia
estadística.

\subsection{Invariantes numéricas}

El build elimina features completamente nulas y verifica ausencia de infinitos. Esto es
diferente de imputar missingness de forma supervisada. Sea $X$ la matriz derivada. La etapa de
datos puede retirar una columna $j$ si
\[
\forall i,\; X_{ij}=\mathrm{NaN},
\]
porque no contiene información en ninguna parcela. En cambio, para una columna parcialmente
observada, la imputación debe estimarse en el conjunto de entrenamiento del fold:
\[
\tilde X_{ij}=
\begin{cases}
X_{ij}, & X_{ij}\text{ observado},\\
\widehat m_j^{(\mathrm{train})}, & X_{ij}\text{ missing}.
\end{cases}
\]
Usar una mediana calculada con el fold de validación no usa $y$, pero en una evaluación
inductiva puede transferir información de distribución. Por disciplina, Checkpoints~02--03
mantuvieron el preprocessing aprendido dentro del pipeline.

\subsection{Representación y provenance}

Cada entrada del manifest registra nombre de columna, fuente, familia y modo. El manifest
funciona como un esquema semántico que permite formar subconjuntos sin depender de regex
frágiles sobre miles de nombres.

Esta estructura fue crítica en Checkpoint~03, donde se construyeron capas agronómicas y
empíricas separadas, y en Checkpoint~04, donde C0, C1 y C4 se resolvieron mediante provenance
en vez de mantener listas manuales de columnas.

\subsection{Por qué no escalar el target}

En discusiones tempranas se consideró multiplicar el rendimiento por 100 para evitar RMSE
numéricamente menor que uno. Esa transformación no mejora el problema. Si
\[
y_i'=c\,y_i,\qquad \hat y_i'=c\,\hat y_i,
\]
entonces
\[
\RMSE(y',\hat y')=|c|\,\RMSE(y,\hat y).
\]
La escala cambia la magnitud reportada, pero no el orden relativo de modelos ni la calidad de
la reconstrucción. Por ello se conserva la unidad natural de toneladas por hectárea.

\subsection{Condición de congelamiento}

Feature Table v1 queda congelada como artefacto reproducible. Las capas posteriores son
\emph{aditivas}: no reescriben las columnas base ni modifican retrospectivamente un build para
hacer que un resultado posterior parezca parte de la representación original. Esta política
hace posible comparar B0, B1, B2, etc. bajo la misma fuente de verdad.
